%!TEX root = ../main.tex
\chapter{Introduction}

The average user's rights and freedoms in the online world have taken a beneficial turn with the introduction of privacy regulations across continents. With the adoption of General Data Protection Regulation (GDPR) in 2016 and its becoming effective two years later, were the first major step in years that granted the users to have control and a say over their information online \cite{gdprhistory}. The California Consumer Privacy Act (CCPA) followed not long after in 2018. 

In contrast to the GDPR's enforcement of \emph{opt-in} mechanisms, the CCPA grants the users the right to \emph{opt-out} of selling and distributing of their personal information to the third parties by businesses which means the consent to not sell or share is asked after a certain amount of information is already collected. The California-based websites that meet the criteria for being categorized as a business are obligated to provide the means to exercise opt-out rights to their visitors \cite[§1798.120 (a)–(b)]{ccpa}. To tackle this, the act requires the business to set up and "provide a clear and conspicuous link on [its]
Internet homepage, titled ‘Do Not Sell My Personal Information’" detailing the user's rights and showing them how to opt out \cite[§1798.135 (a)]{ccpa}. Moreover, the websites must acknowledge and respect the preference signals coming from users' browsers. In the past however these signals lacked a standard. After the group responsible for developing the "do not track" HTTP header was disbanded and the method became obsolete \cite{iapp}, the need for a new universal standard emerged and this paved the way for Global Privacy Control (GPC) to take the reins.

This system granted the users to report to the websites they visit the request to opt out from sharing or selling of their information under the CCPA and the websites that are subject to it are expected to honor it. 




1. Talk about data privacy regulations and website compliance -- mention the paper you are extending.

2. mention how it is very difficult to regulate compliance despite all the precautions like gpc acknowledgement -- Fingerprinting

3. Briefly talk about fingerprinting (details later in background) and talk about project topic and focus.

4. Briefly mention the questions answered and my findings (later)

5. Overview

%This is an example of a numbered chapter. The section number increases automatically and does not need to be stated. To use a unnumbered section (a.k.a. a special section), just use \texttt{\textbackslash{}section*\{Title of unnumbered chapter\}}. The hierarchy is chapter $\rightarrow$ section $\rightarrow$ subsection $\rightarrow$ subsubsection $\rightarrow$ paragraph.


\section{Section}

You guessed it, this is one of those infamous \emph{sections}.

\subsection{Subsections}

This is a numerated \texttt{\textbackslash{}subsection}. There is also a \texttt{\textbackslash{}subsubsection} command. To create an unnumbered subsection, just add * like needed for sections.

\subsubsection{Subsubsection}

And while we're at it, here's a sub-sub-section for you!

\paragraph{A paragraph} is you! But there are no dragons, sorry about that.

